\section*{ПРИЛОЖЕНИЕ В. Надёжность, API-семантика и эксплуатационные сценарии}

\subsection*{В.1. Почему API является частью исследовательского результата}

В инженерных диссертациях API часто воспринимается как сугубо
прикладная деталь. Однако в данной работе API тесно связан с самой
исследовательской гипотезой. Если гипотеза состоит в том, что
состояние должно поддерживаться как snapshot~+~stream, то и интерфейс
сервиса обязан отражать эту модель.

Набор операций \texttt{GetASNInfo}, \texttt{GetPrefixesASTable} и
\texttt{SubscribePrefixASUpdates} не случаен. Он соответствует трём
способам доступа к materialized state \cite{kleppmann2017ddia}:

\begin{itemize}
\item точечное чтение текущего значения;
\item чтение полного среза;
\item подписка на инкрементальные изменения.
\end{itemize}

Именно эта тройка делает внешний контракт сервиса согласованным с его
внутренней моделью.

\subsection*{В.2. Семантика lookup-запроса}

\texttt{GetASNInfo} решает наиболее частую и наиболее простую с точки
зрения клиента задачу: по конкретному IP-адресу определить наиболее
специфичный префикс и связанный с ним ASN. Для downstream-клиента эта
операция должна быть:

\begin{itemize}
\item быстрой;
\item предсказуемой;
\item не требующей знания внутренней структуры таблицы.
\end{itemize}

Внутренне сервис может реализовывать longest prefix match сложным
образом, но наружу он обязан предоставлять простую и стабильную
семантику: «вернуть лучшее известное на текущий момент соответствие».

\subsection*{В.3. Семантика полного снимка}

\texttt{GetPrefixesASTable} на первый взгляд выглядит самой простой
операцией: просто вернуть всю таблицу. На деле именно здесь скрываются
важные вопросы семантики.

Во-первых, «вся таблица» должна быть согласованным snapshot, а не
набором записей, собранных в середине конкурентных изменений.

Во-вторых, клиенту обычно нужна таблица в форме, пригодной для
дальнейшего анализа, а не внутреннее представление дерева.

В-третьих, для крупных таблиц важно понимать стоимость операции и не
считать её «дешёвой по умолчанию» -- особенно в режиме шардированного
ART, где \texttt{Dump} обходит несколько ART-деревьев и выполняет
аллокации.

Это делает \texttt{GetPrefixesASTable} не просто сервисным
convenience method, а отдельным операционным сценарием со своей
ценой и своей архитектурной значимостью.

\subsection*{В.4. Семантика подписки на обновления}

\texttt{SubscribePrefixASUpdates} наиболее полно выражает
streaming-природу системы. С точки зрения клиента этот метод
означает: «после установления подписки я хочу получать все изменения
состояния в порядке их применения локальным сервисом».

Здесь возникают важные эксплуатационные вопросы:

\begin{itemize}
\item нужно ли гарантировать delivery exactly once или at least once
\cite{kreps2011kafka};
\item что считается границей начала подписки;
\item как клиент восстанавливается после разрыва;
\item нужно ли клиенту предварительно получить полный снимок.
\end{itemize}

В текущей архитектуре естественная модель выглядит так:

\begin{enumerate}
\item при необходимости клиент сначала получает полный snapshot;
\item затем открывает stream;
\item stream используется как инкрементальная дельта поверх уже
известного снимка.
\end{enumerate}

Эта модель хорошо согласуется с общей идеей snapshot~+~stream и
легко объясняется downstream-командам, в особенности тем, кто знаком
с шаблонами log-oriented архитектур \cite{kleppmann2017ddia}.

\subsection*{В.5. Поведение при отказах}

Надёжность stream-oriented сервиса определяется не только тем, как он
работает в штатном режиме, но и тем, как он переживает сбои. Для
данной системы критичны как минимум три класса отказов:

\begin{enumerate}
\item временная потеря связи с GoBGP;
\item внутренняя ошибка обработки update-события;
\item разрыв клиентской подписки.
\end{enumerate}

Стратегии работы с этими отказами различны.

При потере связи с GoBGP сервис обязан считать поток потенциально
неполным и инициировать reconnect-and-resync (см. рис.~\ref{fig:lifecycle}).

При внутренней ошибке обработки события сервис должен избегать
продолжения работы в сомнительном состоянии. В production-ориентированной
реализации предпочтительнее пересобрать состояние заново, чем
продолжать эксплуатацию с возможной скрытой рассогласованностью.

При разрыве клиентской подписки ответственность переносится на
клиента: он либо повторно подписывается и принимает модель «сначала
снимок, потом дельта», либо строит собственные механизмы компенсации.

\subsection*{В.6. Эксплуатационный баланс между snapshot и stream}

Одна из важных эксплуатационных идей работы состоит в том, что
snapshot и stream не конкурируют друг с другом. Они не являются двумя
разными режимами, из которых нужно выбрать один. Напротив, они
образуют целостную модель:

\textbf{Snapshot нужен:}

\begin{itemize}
\item при старте;
\item при ресинхронизации;
\item для некоторых внешних клиентов;
\item для валидации и отладки.
\end{itemize}

\textbf{Stream нужен:}

\begin{itemize}
\item для поддержания свежести состояния;
\item для локальных изменений;
\item для downstream-reactive processing;
\item для минимизации общей стоимости обновления.
\end{itemize}

Поэтому архитектурно правильно не противопоставлять snapshot и
stream, а рассматривать их как две фазы жизни одного и того же
состояния \cite{akidau2015dataflow,carbone2015flink}.

\subsection*{В.7. Безопасность и контроль доступа}

Хотя сама работа не фокусируется на вопросах аутентификации и
авторизации, при промышленном внедрении сервис неизбежно становится
частью доверенного внутреннего контура. Это означает, что для
production-эксплуатации полезны:

\begin{itemize}
\item разграничение доступа к полному \texttt{Dump} и к точечному
\texttt{Lookup};
\item отдельный контроль на операции подписки;
\item аудит подключений;
\item rate limiting для особенно дорогих операций.
\end{itemize}

Эти соображения не влияют на основную исследовательскую гипотезу, но
определяют готовность системы к реальному использованию.

\subsection*{В.8. Наблюдаемость и профилирование}

Важный практический аспект разработанного сервиса -- наличие средств
диагностики. Возможность включить \texttt{pprof} и исследовать
поведение отдельных компонентов особенно полезна для систем, где
компромиссы между \texttt{Lookup}, \texttt{Dump} и \texttt{Upsert}
видны только на профилях памяти и CPU. С точки зрения системной
инженерии это означает, что сервис строится не как «чёрный ящик», а
как наблюдаемая подсистема, пригодная для дальнейшей эволюции.
Концептуально это близко к идее BMP \cite{rfc7854}: маршрутизационное
состояние и сервисы, его обрабатывающие, должны быть наблюдаемы как
отдельный поток.

\subsection*{В.9. Потенциальное расширение API}

В будущем API сервиса может быть расширен в нескольких направлениях:

\begin{itemize}
\item выдача метаданных о свежести локального состояния (timestamp
последнего обновления, lag относительно источника);
\item получение дифференциального snapshot за интервал времени;
\item выдача статистики по flapping prefixes (с использованием идей
RFC~2439 \cite{rfc2439});
\item запросы по ASN с получением всех известных префиксов (обратное
отношение);
\item диагностические методы для состояния sync/watch-loop;
\item поддержка фильтрации по \texttt{<AFI,SAFI>}
\cite{rfc4760} в потоке подписки.
\end{itemize}

Все эти направления логически следуют из уже реализованного ядра и
усиливают ценность сервиса как общей платформы маршрутизационной
аналитики.

\subsection*{В.10. Итог по эксплуатационной части}

Эксплуатационная ценность разработанного решения состоит не только в
том, что оно «быстро обновляет дерево». Она состоит в том, что
построен сервис с корректной state model, ясным API и внятным
поведением при отказах. Именно сочетание алгоритмического ядра и
системной оболочки делает результат диссертационным, а не просто
локальной оптимизацией структуры данных.

%==============================================================
